% --- Archon physics formalization source begin ---
% archon:physics
% archon:covers PhyXMiniProblems/problem_phyx_mini_0412.lean
% archon:source-report reports/phyx_mini/problem_phyx_mini_0412.source.json

\chapter{Physics problem phyx\_mini\_0412}
\label{ch:PhyXMiniProblems_problem_phyx_mini_0412}

\paragraph{Problem source.}
\#\# Physical scenario

A monatomic gas follows the process \$1 \textbackslash{}rightarrow 2 \textbackslash{}rightarrow 3\$ shown in figure.

\#\# Question

How much heat is needed for process \$2 \textbackslash{}rightarrow 3\$?

\#\# Answer choices

- A: A: 91 J

- B: B: 0 J

- C: C: -91 J

- D: D: -150 J

\#\# Figure caption (auxiliary; use the image as primary evidence)

The image is a pressure-volume (p-V) diagram depicting a thermodynamic process. The axes are labeled "p (atm)" for pressure in atmospheres (y-axis) and "V (cm³)" for volume in cubic centimeters (x-axis).

There are three marked points: 1, 2, and 3. 

- Point 1 is at approximately 3 atm and 100 cm³.
- Point 2 is at approximately 3 atm and 300 cm³.
- Point 3 is at approximately 1 atm and 300 cm³.

The process from point 1 to point 2 is an isobaric expansion (constant pressure) shown by a horizontal line with an arrow pointing right. The process from point 2 to point 3 is an isochoric cooling (constant volume) shown by a vertical line with an arrow pointing downward.

A red dashed curve labeled "100°C isotherm" runs between point 1 and point 3, representing an isothermal process at 100°C that the system does not follow. The curve shows the relationship between pressure and volume at constant temperature.

\#\# Dataset metadata

Laws of Thermodynamics; Physical Model Grounding Reasoning, Multi-Formula Reasoning

\paragraph{Recorded answer/context.}
C: C: -91 J

\paragraph{Figure/image path.}
/root/proposal\_for\_physic/science-mango/phyx\_mini\_run/phyx\_data/test\_image/412.png

\paragraph{Formalization target.}
create a compiling Lean file with sorry bodies at `PhyXMiniProblems/problem\_phyx\_mini\_0412.lean`. The Lean declarations must preserve the physical quantities, dimensions or dimensional roles, figure labels, governing-law hypotheses, and final relation expressed by this problem.
Use Mathlib/Physlib names found through LeanExplore where available. If a physics API is missing, introduce faithful local abstractions rather than scalar placeholder aliases.

\begin{theorem}[Physics formalization target]
\label{thm:physics:phyx_mini_0412:target}
\lean{PhyXMiniProblems.ProblemPhyXMini0412.heatForTwoToThree_matches_recordedAnswerC}
\uses{def:physics:phyx-mini-0412:phyxminiproblems-problemphyxmini0412-energyinjoules, def:physics:phyx-mini-0412:phyxminiproblems-problemphyxmini0412-monatomicgasprocesssetup, def:physics:phyx-mini-0412:phyxminiproblems-problemphyxmini0412-matchesproblemandprimaryfigure, def:physics:phyx-mini-0412:phyxminiproblems-problemphyxmini0412-hasphysicalstatecoordinates, def:physics:phyx-mini-0412:phyxminiproblems-problemphyxmini0412-obeysmonatomicidealgasthermodynamics, def:physics:phyx-mini-0412:phyxminiproblems-problemphyxmini0412-displayedheatinjoules, def:physics:phyx-mini-0412:phyxminiproblems-problemphyxmini0412-recordedanswerchoice, def:physics:phyx-mini-0412:phyxminiproblems-problemphyxmini0412-roundstonearestwholejoule}
The assigned autoformalize agent should translate this physics problem into Lean declarations in the covered file, with theorem and lemma proof bodies written as `by sorry`.
\end{theorem}
\begin{proof}
This is an autoformalization task, not a proof task. Produce statements that can later be proved by the physics prover without weakening the physical model.
\end{proof}

% --- Archon declaration topology begin ---
\section{Declaration topology}
\begin{definition}[Volume Quantity]
  \label{def:physics:phyx-mini-0412:phyxminiproblems-problemphyxmini0412-volumequantity}
  \lean{PhyXMiniProblems.ProblemPhyXMini0412.VolumeQuantity}
  A nonnegative physical volume, carrying dimension length³
\end{definition}
\begin{proof}
  This is the declared model definition of Volume Quantity.
\end{proof}
\begin{definition}[Pressure Quantity]
  \label{def:physics:phyx-mini-0412:phyxminiproblems-problemphyxmini0412-pressurequantity}
  \lean{PhyXMiniProblems.ProblemPhyXMini0412.PressureQuantity}
  Thermodynamic pressure, represented by Physlib's dimensionful pressure type
\end{definition}
\begin{proof}
  This is the declared model definition of Pressure Quantity.
\end{proof}
\begin{definition}[Temperature Quantity]
  \label{def:physics:phyx-mini-0412:phyxminiproblems-problemphyxmini0412-temperaturequantity}
  \lean{PhyXMiniProblems.ProblemPhyXMini0412.TemperatureQuantity}
  A nonnegative absolute temperature, carrying the temperature dimension
\end{definition}
\begin{proof}
  This is the declared model definition of Temperature Quantity.
\end{proof}
\begin{definition}[Energy Quantity]
  \label{def:physics:phyx-mini-0412:phyxminiproblems-problemphyxmini0412-energyquantity}
  \lean{PhyXMiniProblems.ProblemPhyXMini0412.EnergyQuantity}
  Signed heat, work, or internal energy, represented by dimensionful energy
\end{definition}
\begin{proof}
  This is the declared model definition of Energy Quantity.
\end{proof}
\begin{definition}[volume In Cubic Meters]
  \label{def:physics:phyx-mini-0412:phyxminiproblems-problemphyxmini0412-volumeincubicmeters}
  \lean{PhyXMiniProblems.ProblemPhyXMini0412.volumeInCubicMeters}
  \uses{def:physics:phyx-mini-0412:phyxminiproblems-problemphyxmini0412-volumequantity}
  SI cubic-metre readout of a physical volume
\end{definition}
\begin{proof}
  This is the declared model definition of volume In Cubic Meters.
\end{proof}
\begin{definition}[volume In Cubic Centimeters]
  \label{def:physics:phyx-mini-0412:phyxminiproblems-problemphyxmini0412-volumeincubiccentimeters}
  \lean{PhyXMiniProblems.ProblemPhyXMini0412.volumeInCubicCentimeters}
  \uses{def:physics:phyx-mini-0412:phyxminiproblems-problemphyxmini0412-volumequantity, def:physics:phyx-mini-0412:phyxminiproblems-problemphyxmini0412-volumeincubicmeters}
  Cubic-centimetre readout used by the horizontal axis of the figure
\end{definition}
\begin{proof}
  This is the declared model definition of volume In Cubic Centimeters.
\end{proof}
\begin{definition}[pressure In Pascals]
  \label{def:physics:phyx-mini-0412:phyxminiproblems-problemphyxmini0412-pressureinpascals}
  \lean{PhyXMiniProblems.ProblemPhyXMini0412.pressureInPascals}
  \uses{def:physics:phyx-mini-0412:phyxminiproblems-problemphyxmini0412-pressurequantity}
  Pascal readout of a physical pressure
\end{definition}
\begin{proof}
  This is the declared model definition of pressure In Pascals.
\end{proof}
\begin{definition}[pressure In Standard Atmospheres]
  \label{def:physics:phyx-mini-0412:phyxminiproblems-problemphyxmini0412-pressureinstandardatmospheres}
  \lean{PhyXMiniProblems.ProblemPhyXMini0412.pressureInStandardAtmospheres}
  \uses{def:physics:phyx-mini-0412:phyxminiproblems-problemphyxmini0412-pressurequantity}
  Standard-atmosphere readout used by the vertical axis of the figure
\end{definition}
\begin{proof}
  This is the declared model definition of pressure In Standard Atmospheres.
\end{proof}
\begin{definition}[temperature In Kelvins]
  \label{def:physics:phyx-mini-0412:phyxminiproblems-problemphyxmini0412-temperatureinkelvins}
  \lean{PhyXMiniProblems.ProblemPhyXMini0412.temperatureInKelvins}
  \uses{def:physics:phyx-mini-0412:phyxminiproblems-problemphyxmini0412-temperaturequantity}
  SI kelvin readout of an absolute temperature
\end{definition}
\begin{proof}
  This is the declared model definition of temperature In Kelvins.
\end{proof}
\begin{definition}[temperature In Degrees Celsius]
  \label{def:physics:phyx-mini-0412:phyxminiproblems-problemphyxmini0412-temperatureindegreescelsius}
  \lean{PhyXMiniProblems.ProblemPhyXMini0412.temperatureInDegreesCelsius}
  \uses{def:physics:phyx-mini-0412:phyxminiproblems-problemphyxmini0412-temperaturequantity, def:physics:phyx-mini-0412:phyxminiproblems-problemphyxmini0412-temperatureinkelvins}
  Degrees-Celsius readout used by the dashed isotherm annotation
\end{definition}
\begin{proof}
  This is the declared model definition of temperature In Degrees Celsius.
\end{proof}
\begin{definition}[energy In Joules]
  \label{def:physics:phyx-mini-0412:phyxminiproblems-problemphyxmini0412-energyinjoules}
  \lean{PhyXMiniProblems.ProblemPhyXMini0412.energyInJoules}
  \uses{def:physics:phyx-mini-0412:phyxminiproblems-problemphyxmini0412-energyquantity}
  Joule readout of a signed dimensionful energy quantity
\end{definition}
\begin{proof}
  This is the declared model definition of energy In Joules.
\end{proof}
\begin{definition}[atmosphere Cubic Centimeter In Joules]
  \label{def:physics:phyx-mini-0412:phyxminiproblems-problemphyxmini0412-atmospherecubiccentimeterinjoules}
  \lean{PhyXMiniProblems.ProblemPhyXMini0412.atmosphereCubicCentimeterInJoules}
  One standard atmosphere times one cubic centimetre, expressed in joules. This is the unit conversion 101325 Pa × 10⁻⁶ m³ = 0.101325 J
\end{definition}
\begin{proof}
  This is the declared model definition of atmosphere Cubic Centimeter In Joules.
\end{proof}
\begin{definition}[Gas Atomicity]
  \label{def:physics:phyx-mini-0412:phyxminiproblems-problemphyxmini0412-gasatomicity}
  \lean{PhyXMiniProblems.ProblemPhyXMini0412.GasAtomicity}
  Atomicity of the gas specified by the problem
\end{definition}
\begin{proof}
  This is the declared model definition of Gas Atomicity.
\end{proof}
\begin{definition}[Gas Model]
  \label{def:physics:phyx-mini-0412:phyxminiproblems-problemphyxmini0412-gasmodel}
  \lean{PhyXMiniProblems.ProblemPhyXMini0412.GasModel}
  Equation-of-state model used for the gas
\end{definition}
\begin{proof}
  This is the declared model definition of Gas Model.
\end{proof}
\begin{definition}[Gas Sample]
  \label{def:physics:phyx-mini-0412:phyxminiproblems-problemphyxmini0412-gassample}
  \lean{PhyXMiniProblems.ProblemPhyXMini0412.GasSample}
  \uses{def:physics:phyx-mini-0412:phyxminiproblems-problemphyxmini0412-gasatomicity, def:physics:phyx-mini-0412:phyxminiproblems-problemphyxmini0412-gasmodel}
  A physical gas sample, distinguished from any of its scalar readouts
\end{definition}
\begin{proof}
  This is the declared model definition of Gas Sample.
\end{proof}
\begin{definition}[State Label]
  \label{def:physics:phyx-mini-0412:phyxminiproblems-problemphyxmini0412-statelabel}
  \lean{PhyXMiniProblems.ProblemPhyXMini0412.StateLabel}
  The three state labels printed in the pressure-volume diagram
\end{definition}
\begin{proof}
  This is the declared model definition of State Label.
\end{proof}
\begin{definition}[Process Leg]
  \label{def:physics:phyx-mini-0412:phyxminiproblems-problemphyxmini0412-processleg}
  \lean{PhyXMiniProblems.ProblemPhyXMini0412.ProcessLeg}
  The two directed process legs followed by the gas
\end{definition}
\begin{proof}
  This is the declared model definition of Process Leg.
\end{proof}
\begin{definition}[Process Kind]
  \label{def:physics:phyx-mini-0412:phyxminiproblems-problemphyxmini0412-processkind}
  \lean{PhyXMiniProblems.ProblemPhyXMini0412.ProcessKind}
  Thermodynamic classification of a depicted process leg
\end{definition}
\begin{proof}
  This is the declared model definition of Process Kind.
\end{proof}
\begin{definition}[Process Direction]
  \label{def:physics:phyx-mini-0412:phyxminiproblems-problemphyxmini0412-processdirection}
  \lean{PhyXMiniProblems.ProblemPhyXMini0412.ProcessDirection}
  Qualitative change indicated by a directed process arrow
\end{definition}
\begin{proof}
  This is the declared model definition of Process Direction.
\end{proof}
\begin{definition}[Thermodynamic State]
  \label{def:physics:phyx-mini-0412:phyxminiproblems-problemphyxmini0412-thermodynamicstate}
  \lean{PhyXMiniProblems.ProblemPhyXMini0412.ThermodynamicState}
  \uses{def:physics:phyx-mini-0412:phyxminiproblems-problemphyxmini0412-volumequantity, def:physics:phyx-mini-0412:phyxminiproblems-problemphyxmini0412-pressurequantity, def:physics:phyx-mini-0412:phyxminiproblems-problemphyxmini0412-temperaturequantity, def:physics:phyx-mini-0412:phyxminiproblems-problemphyxmini0412-energyquantity}
  Pressure, volume, temperature, and internal energy at equilibrium
\end{definition}
\begin{proof}
  This is the declared model definition of Thermodynamic State.
\end{proof}
\begin{definition}[Thermodynamic Process]
  \label{def:physics:phyx-mini-0412:phyxminiproblems-problemphyxmini0412-thermodynamicprocess}
  \lean{PhyXMiniProblems.ProblemPhyXMini0412.ThermodynamicProcess}
  \uses{def:physics:phyx-mini-0412:phyxminiproblems-problemphyxmini0412-energyquantity, def:physics:phyx-mini-0412:phyxminiproblems-problemphyxmini0412-statelabel, def:physics:phyx-mini-0412:phyxminiproblems-problemphyxmini0412-processkind, def:physics:phyx-mini-0412:phyxminiproblems-problemphyxmini0412-processdirection}
  A directed thermodynamic process. Heat is positive when transferred into the gas, and boundary work is positive when done by the gas. Neither energy field is defined from an answer choice
\end{definition}
\begin{proof}
  This is the declared model definition of Thermodynamic Process.
\end{proof}
\begin{definition}[Figure Axis]
  \label{def:physics:phyx-mini-0412:phyxminiproblems-problemphyxmini0412-figureaxis}
  \lean{PhyXMiniProblems.ProblemPhyXMini0412.FigureAxis}
  The two Cartesian axes in the primary pressure-volume diagram
\end{definition}
\begin{proof}
  This is the declared model definition of Figure Axis.
\end{proof}
\begin{definition}[Axis Quantity]
  \label{def:physics:phyx-mini-0412:phyxminiproblems-problemphyxmini0412-axisquantity}
  \lean{PhyXMiniProblems.ProblemPhyXMini0412.AxisQuantity}
  Physical role assigned to a figure axis
\end{definition}
\begin{proof}
  This is the declared model definition of Axis Quantity.
\end{proof}
\begin{definition}[Axis Display Unit]
  \label{def:physics:phyx-mini-0412:phyxminiproblems-problemphyxmini0412-axisdisplayunit}
  \lean{PhyXMiniProblems.ProblemPhyXMini0412.AxisDisplayUnit}
  Unit text printed beside a figure axis
\end{definition}
\begin{proof}
  This is the declared model definition of Axis Display Unit.
\end{proof}
\begin{definition}[Axis Symbol]
  \label{def:physics:phyx-mini-0412:phyxminiproblems-problemphyxmini0412-axissymbol}
  \lean{PhyXMiniProblems.ProblemPhyXMini0412.AxisSymbol}
  Literal physical-quantity symbol printed beside a figure axis
\end{definition}
\begin{proof}
  This is the declared model definition of Axis Symbol.
\end{proof}
\begin{definition}[Path Geometry]
  \label{def:physics:phyx-mini-0412:phyxminiproblems-problemphyxmini0412-pathgeometry}
  \lean{PhyXMiniProblems.ProblemPhyXMini0412.PathGeometry}
  Geometric appearance of a process leg in the p-V plane
\end{definition}
\begin{proof}
  This is the declared model definition of Path Geometry.
\end{proof}
\begin{definition}[Figure Line Style]
  \label{def:physics:phyx-mini-0412:phyxminiproblems-problemphyxmini0412-figurelinestyle}
  \lean{PhyXMiniProblems.ProblemPhyXMini0412.FigureLineStyle}
  Visible line style used for a feature of the diagram
\end{definition}
\begin{proof}
  This is the declared model definition of Figure Line Style.
\end{proof}
\begin{definition}[Pressure Volume Figure]
  \label{def:physics:phyx-mini-0412:phyxminiproblems-problemphyxmini0412-pressurevolumefigure}
  \lean{PhyXMiniProblems.ProblemPhyXMini0412.PressureVolumeFigure}
  \uses{def:physics:phyx-mini-0412:phyxminiproblems-problemphyxmini0412-volumequantity, def:physics:phyx-mini-0412:phyxminiproblems-problemphyxmini0412-pressurequantity, def:physics:phyx-mini-0412:phyxminiproblems-problemphyxmini0412-statelabel, def:physics:phyx-mini-0412:phyxminiproblems-problemphyxmini0412-processleg, def:physics:phyx-mini-0412:phyxminiproblems-problemphyxmini0412-figureaxis, def:physics:phyx-mini-0412:phyxminiproblems-problemphyxmini0412-axisquantity, def:physics:phyx-mini-0412:phyxminiproblems-problemphyxmini0412-axisdisplayunit, def:physics:phyx-mini-0412:phyxminiproblems-problemphyxmini0412-axissymbol, def:physics:phyx-mini-0412:phyxminiproblems-problemphyxmini0412-pathgeometry, def:physics:phyx-mini-0412:phyxminiproblems-problemphyxmini0412-figurelinestyle}
  Structured transcription of image 412.png. The dashed isotherm is retained as a figure feature separate from the two solid paths actually followed
\end{definition}
\begin{proof}
  This is the declared model definition of Pressure Volume Figure.
\end{proof}
\begin{definition}[Monatomic Gas Process Setup]
  \label{def:physics:phyx-mini-0412:phyxminiproblems-problemphyxmini0412-monatomicgasprocesssetup}
  \lean{PhyXMiniProblems.ProblemPhyXMini0412.MonatomicGasProcessSetup}
  \uses{def:physics:phyx-mini-0412:phyxminiproblems-problemphyxmini0412-gassample, def:physics:phyx-mini-0412:phyxminiproblems-problemphyxmini0412-statelabel, def:physics:phyx-mini-0412:phyxminiproblems-problemphyxmini0412-processleg, def:physics:phyx-mini-0412:phyxminiproblems-problemphyxmini0412-thermodynamicstate, def:physics:phyx-mini-0412:phyxminiproblems-problemphyxmini0412-thermodynamicprocess, def:physics:phyx-mini-0412:phyxminiproblems-problemphyxmini0412-pressurevolumefigure}
  The gas, its three states, the two process legs, and the supplied figure
\end{definition}
\begin{proof}
  This is the declared model definition of Monatomic Gas Process Setup.
\end{proof}
\begin{definition}[Matches Problem And Primary Figure]
  \label{def:physics:phyx-mini-0412:phyxminiproblems-problemphyxmini0412-matchesproblemandprimaryfigure}
  \lean{PhyXMiniProblems.ProblemPhyXMini0412.MatchesProblemAndPrimaryFigure}
  \uses{def:physics:phyx-mini-0412:phyxminiproblems-problemphyxmini0412-volumeincubiccentimeters, def:physics:phyx-mini-0412:phyxminiproblems-problemphyxmini0412-pressureinstandardatmospheres, def:physics:phyx-mini-0412:phyxminiproblems-problemphyxmini0412-temperatureindegreescelsius, def:physics:phyx-mini-0412:phyxminiproblems-problemphyxmini0412-statelabel, def:physics:phyx-mini-0412:phyxminiproblems-problemphyxmini0412-processleg, def:physics:phyx-mini-0412:phyxminiproblems-problemphyxmini0412-monatomicgasprocesssetup}
  Exact transcription of the prose and primary bitmap. The figure fixes state coordinates, process classifications, directions, and the dashed 100°C isotherm. It supplies no numerical heat, work, or internal-energy value
\end{definition}
\begin{proof}
  This is the declared model definition of Matches Problem And Primary Figure.
\end{proof}
\begin{definition}[Has Physical State Coordinates]
  \label{def:physics:phyx-mini-0412:phyxminiproblems-problemphyxmini0412-hasphysicalstatecoordinates}
  \lean{PhyXMiniProblems.ProblemPhyXMini0412.HasPhysicalStateCoordinates}
  \uses{def:physics:phyx-mini-0412:phyxminiproblems-problemphyxmini0412-volumeincubicmeters, def:physics:phyx-mini-0412:phyxminiproblems-problemphyxmini0412-pressureinpascals, def:physics:phyx-mini-0412:phyxminiproblems-problemphyxmini0412-temperatureinkelvins, def:physics:phyx-mini-0412:phyxminiproblems-problemphyxmini0412-statelabel, def:physics:phyx-mini-0412:phyxminiproblems-problemphyxmini0412-monatomicgasprocesssetup}
  Positivity conditions selecting physically meaningful state coordinates
\end{definition}
\begin{proof}
  This is the declared model definition of Has Physical State Coordinates.
\end{proof}
\begin{definition}[Obeys Monatomic Ideal Gas Thermodynamics]
  \label{def:physics:phyx-mini-0412:phyxminiproblems-problemphyxmini0412-obeysmonatomicidealgasthermodynamics}
  \lean{PhyXMiniProblems.ProblemPhyXMini0412.ObeysMonatomicIdealGasThermodynamics}
  \uses{def:physics:phyx-mini-0412:phyxminiproblems-problemphyxmini0412-volumeincubiccentimeters, def:physics:phyx-mini-0412:phyxminiproblems-problemphyxmini0412-pressureinstandardatmospheres, def:physics:phyx-mini-0412:phyxminiproblems-problemphyxmini0412-energyinjoules, def:physics:phyx-mini-0412:phyxminiproblems-problemphyxmini0412-atmospherecubiccentimeterinjoules, def:physics:phyx-mini-0412:phyxminiproblems-problemphyxmini0412-statelabel, def:physics:phyx-mini-0412:phyxminiproblems-problemphyxmini0412-processleg, def:physics:phyx-mini-0412:phyxminiproblems-problemphyxmini0412-monatomicgasprocesssetup}
  The laws used for a monatomic ideal gas, with explicit unit readouts: U = (3/2) pV at each equilibrium state; Q\_in = ΔU + W\_by for each directed process; and an isochoric leg has zero boundary work. All three relations are generic over the setup's states or legs and contain no problem-specific heat value or answer choice
\end{definition}
\begin{proof}
  This is the declared model definition of Obeys Monatomic Ideal Gas Thermodynamics.
\end{proof}
\begin{definition}[Answer Choice]
  \label{def:physics:phyx-mini-0412:phyxminiproblems-problemphyxmini0412-answerchoice}
  \lean{PhyXMiniProblems.ProblemPhyXMini0412.AnswerChoice}
  The four labels printed beside the multiple-choice answers
\end{definition}
\begin{proof}
  This is the declared model definition of Answer Choice.
\end{proof}
\begin{definition}[displayed Heat In Joules]
  \label{def:physics:phyx-mini-0412:phyxminiproblems-problemphyxmini0412-displayedheatinjoules}
  \lean{PhyXMiniProblems.ProblemPhyXMini0412.displayedHeatInJoules}
  \uses{def:physics:phyx-mini-0412:phyxminiproblems-problemphyxmini0412-answerchoice}
  Heat readout in joules printed beside each answer label
\end{definition}
\begin{proof}
  This is the declared model definition of displayed Heat In Joules.
\end{proof}
\begin{definition}[recorded Answer Choice]
  \label{def:physics:phyx-mini-0412:phyxminiproblems-problemphyxmini0412-recordedanswerchoice}
  \lean{PhyXMiniProblems.ProblemPhyXMini0412.recordedAnswerChoice}
  \uses{def:physics:phyx-mini-0412:phyxminiproblems-problemphyxmini0412-answerchoice}
  The answer label recorded in the dataset metadata
\end{definition}
\begin{proof}
  This is the declared model definition of recorded Answer Choice.
\end{proof}
\begin{definition}[Rounds To Nearest Whole Joule]
  \label{def:physics:phyx-mini-0412:phyxminiproblems-problemphyxmini0412-roundstonearestwholejoule}
  \lean{PhyXMiniProblems.ProblemPhyXMini0412.RoundsToNearestWholeJoule}
  \uses{def:physics:phyx-mini-0412:phyxminiproblems-problemphyxmini0412-energyquantity, def:physics:phyx-mini-0412:phyxminiproblems-problemphyxmini0412-energyinjoules}
  A unit-labelled heat readout rounds to a displayed whole-joule value when it lies strictly within half a joule of that value. This generic display rule contains no problem-specific answer
\end{definition}
\begin{proof}
  This is the declared model definition of Rounds To Nearest Whole Joule.
\end{proof}
\begin{lemma}[heat For Two To Three from monatomic first law]
  \label{lem:physics:phyx-mini-0412:phyxminiproblems-problemphyxmini0412-heatfortwotothree-from-monatomic-first-law}
  \lean{PhyXMiniProblems.ProblemPhyXMini0412.heatForTwoToThree_from_monatomic_first_law}
  \uses{def:physics:phyx-mini-0412:phyxminiproblems-problemphyxmini0412-energyinjoules, def:physics:phyx-mini-0412:phyxminiproblems-problemphyxmini0412-atmospherecubiccentimeterinjoules, def:physics:phyx-mini-0412:phyxminiproblems-problemphyxmini0412-monatomicgasprocesssetup, def:physics:phyx-mini-0412:phyxminiproblems-problemphyxmini0412-matchesproblemandprimaryfigure, def:physics:phyx-mini-0412:phyxminiproblems-problemphyxmini0412-obeysmonatomicidealgasthermodynamics}
  The governing laws and the state-2/state-3 figure data give Q₂₃ = (3/2)(1 atm·300 cm³ - 3 atm·300 cm³). This intermediate heat relation is a conclusion, not a premise field
\end{lemma}
\begin{proof}
  The conclusion follows from the governing relations and figure readouts named in the statement of heat For Two To Three from monatomic first law.
\end{proof}
% --- Archon declaration topology end ---
% --- Archon physics formalization source end ---
